% Pillar-7: signal-starvation theory (T1-T3), each theorem paired with its check
\section{A Signal-Starvation Model of Binary-Reward GRPO}
\label{sec:p7-theory}

\subsection{Setup and the usable-signal functional}
\label{sec:p7-theory-setup}
Fix a prompt $x$ with latent per-rollout success probability $p \in [0,1]$
under the current policy, and let a GRPO group consist of $G$ conditionally
i.i.d.\ binary rewards $R_1, \dots, R_G \sim \mathrm{Bernoulli}(p)$, so the
success count is $K = \sum_i R_i \sim \mathrm{Binomial}(G, p)$. The group is
\emph{degenerate}---zero within-group variance, zero centered advantage, zero
gradient---iff $K \in \{0, G\}$. Define
\begin{equation}
  h_G(p) \;=\; p^G + (1-p)^G,
  \label{eq:p7-hG}
\end{equation}
the probability of degeneracy, and model the \emph{usable learning signal}
per prompt as
\begin{equation}
  S(p, G) \;=\; p(1-p)\,\bigl(1 - h_G(p)\bigr).
  \label{eq:p7-S}
\end{equation}
The first factor is the contrast magnitude a non-degenerate group can carry
(the Bernoulli variance, which also bounds the mean absolute advantage;
Appendix~\ref{sec:p7-appendix}); the second is the probability that any
signal survives group centering at all. Both factors vanish at $p \in
\{0,1\}$ and peak at $p = \nicefrac{1}{2}$: GRPO learns fastest exactly where
the policy is most uncertain, and starves at both ends. Three consequences
follow; proofs are deferred to Appendix~\ref{sec:p7-appendix}, and each
theorem is paired here with the measurement that could have falsified it.

\subsection{T1: the expected ZVF indicator is $p^G + (1-p)^G$}
\label{sec:p7-t1}
\textbf{Theorem 1.} \emph{Under the binomial model, the expected value of the
per-group ZVF indicator is exactly $h_G(p) = p^G + (1-p)^G$; aggregated over
a prompt population $\mathcal{D}$, $\E[\mathrm{ZVF}] =
\E_{x \sim \mathcal{D}}[\,p_x^G + (1-p_x)^G\,]$.}

\emph{Empirical check.} Binning the per-group reward tensors of the GSM8K
runs (the only runs that log full per-group tensors; $G = 8$, 3 seeds) by
empirical success rate $\hat p_x$ reproduces the predicted profile---by
construction at the endpoints, and empirically in the interior collapse: the
ZVF indicator averages
$1.000$ at $\hat p_x \in \{0, 1\}$ and $0.000$ at every interior bin
$\hat p_x \in \{0.125, \dots, 0.875\}$
(\texttt{experiments/results/pcd\_vs\_zvf\_shape.tsv}). At the run level, the
group-size sweep of Table~\ref{tab:p7-gsweep} logs the closed-form prediction
$h_G$ evaluated at the run's mean success rate next to the measured mean ZVF:
the two co-vary monotonically in $G$, though measured ZVF sits above the
prediction at every $G$ because the binomial model evaluates $h_G$ at the
\emph{mean} $p$ while the true prompt population mixes mastered and hopeless
prompts whose $h_G \approx 1$ (Jensen gap; this is the U-shape of
Section~\ref{sec:p7-ushape} acting within a single run).

\subsection{T2: larger $G$ lowers ZVF}
\label{sec:p7-t2}
\textbf{Theorem 2.} \emph{For every $p \in (0,1)$, $h_G(p)$ is strictly
decreasing in $G$; hence expected ZVF falls monotonically as group size
grows, with exponential decay rate $-\ln\max(p, 1-p)$.}

\emph{Empirical check (two independent grids).} First, a model~$\times$~$G$
grid from the 368-run audit (Table~\ref{tab:p7-modelG}): for every model with
an interior success rate, measured ZVF falls monotonically in $G$---most
cleanly for Qwen3-32B ($0.33 \to 0.18 \to 0.08$) and Llama-3.1-8B
($0.76 \to 0.65 \to 0.47$). The two boundary cases are exactly the ones the
theory flags: Llama-3.2-1B sits at $p \approx 0$ where $h_G \approx 1$
regardless of $G$ ($0.98 \to 0.97$), and Qwen3-4B-Instruct sits near
saturation where the decay rate $-\ln\max(p,1-p)$ is tiny and seed noise
dominates ($0.89 \to 0.84 \to 0.91$, the single non-monotone cell). Second,
the repo-internal sweep at $G \in \{2,4,8,16\}$ (3 seeds each,
Table~\ref{tab:p7-gsweep}) shows the same monotone decline, mean ZVF
$0.838 \to 0.631$, while held-out accuracy stays flat ($\approx 0.98$
throughout)---raising $G$ buys signal, not (on this near-saturated task)
outcome, a point the controller design of Section~\ref{sec:p7-rules} must
respect.

\begin{table}[t]
\centering
\caption{T2 check 1: mean ZVF by model and group size $G$ (368-run audit,
W\&B project \texttt{zvf-audit}). ZVF falls with $G$ wherever the success
rate is interior; the boundary rows behave as $h_G$ predicts (see text).
``---'' = cell not run.}
\label{tab:p7-modelG}
\begin{tabular}{lccc}
\toprule
Model & $G{=}4$ & $G{=}8$ & $G{=}16$ \\
\midrule
Llama-3.2-1B        & 0.98 & 0.97 & --- \\
Qwen3-4B-Instruct   & 0.89 & 0.84 & 0.91 \\
Llama-3.1-8B        & 0.76 & 0.65 & 0.47 \\
Qwen3-8B            & 0.36 & 0.27 & 0.21 \\
Qwen3-32B           & 0.33 & 0.18 & 0.08 \\
\bottomrule
\end{tabular}
\end{table}

\begin{table}[t]
\centering
\caption{T2 check 2: repo-internal group-size sweep (3 seeds per $G$;
\texttt{experiments/results/groupsize\_zvf\_sweep.tsv}). Measured mean ZVF
falls monotonically in $G$ and co-varies with the closed-form prediction
$h_G(\bar p)$ evaluated at the run's mean training reward $\bar p$; held-out
accuracy is flat because the task is near-saturated.}
\label{tab:p7-gsweep}
\begin{tabular}{ccccc}
\toprule
$G$ & Mean ZVF & $h_G(\bar p)$ & $\bar p$ (train reward) & Held-out acc.\ (SE) \\
\midrule
2  & 0.838 & 0.731 & 0.840 & 0.982 (0.004) \\
4  & 0.764 & 0.554 & 0.863 & 0.988 (0.002) \\
8  & 0.691 & 0.326 & 0.869 & 0.990 (0.003) \\
16 & 0.631 & 0.114 & 0.873 & 0.978 (0.006) \\
\bottomrule
\end{tabular}
\end{table}

\subsection{T3: gradient magnitude tracks $p(1-p)$}
\label{sec:p7-t3}
\textbf{Theorem 3.} \emph{With unnormalized group-centered advantages
$A_i = R_i - \hat p$, the mean absolute advantage in a group is exactly
$2\hat p(1 - \hat p)$; hence the advantage mass driving the policy gradient
scales with the within-group contrast $\hat p(1-\hat p)$, up to the
policy-dependent score-function factor.}

\emph{Empirical check.} Theorem~3 is the load-bearing assumption behind the
first factor of Eq.~\eqref{eq:p7-S}, and it cannot be checked on the closed
stack, which does not expose per-group gradients. Section~\ref{sec:p7-e1}
reports a direct open-backward-pass measurement:
$\mathrm{corr}(\|\nabla\|, \hat p(1-\hat p)) = +0.71$ on a 0.5B model with
synthetic arithmetic---consistent in sign and rank with T3, and scoped
strictly as directional evidence.
